\subsubsection{Translação da base}
\label{sec:dinAcoplada_translacao_base}

Do ponto de vista físico, o movimento do centro de massa da espaçonave
robótica é governado pelas leis de Newton e pelas equações do movimento
orbital.
Em um modelo completo de base flutuante, os graus de liberdade de
translação do centro de massa da espaçonave aparecem no vetor de
coordenadas generalizadas, acoplados à atitude da base e ao movimento
das juntas do \gls{mr}.
Portanto, deslocamentos internos de massa, causados pelo
movimento dos elos, modificam a posição do centro de massa e podem, em
princípio, perturbar o movimento orbital \cite{wie2008space,featherstone2008}.

Neste trabalho, adota-se uma simplificação onde
a translação do centro de massa da espaçonave é descrita pelas equações \gls{hcw} no referencial \gls{lvlh},
apresentado na Seção~\ref{sec:hillcw}, enquanto a formulação
lagrangiana da dinâmica acoplada considera apenas as variáveis de
atitude da base e as juntas do \gls{mr}.
Assim, o vetor de coordenadas generalizadas associado à dinâmica
espaçonave--manipulador é dado por \eqref{eq:q_generalizado_base_mr}, sem
incluir explicitamente coordenadas de posição da base no modelo
lagrangiano.

Em resumo, as equações do movimento orbital relativo fornecem a cada
instante a posição e a velocidade do centro de massa da espaçonave no
referencial \gls{lvlh}, que são tratadas como dados de entrada para o
modelo de atitude e juntas.
As reações devidas ao movimento interno do \gls{mr} sobre a
órbita do centro de massa não são realimentadas nas equações de
translação.
Por outro lado, os acoplamentos rotacionais entre base e manipulador
são mantidos explicitamente na matriz de inércia e nos termos de
Coriolis em \eqref{eq:dinamica_geral}, de modo que o efeito dos torques
das juntas sobre a atitude da espaçonave é preservado na dinâmica
acoplada.

Essa hipótese de desacoplamento entre a translação do centro de massa e
as reações internas é adequada para analisar, de forma isolada, o
comportamento de atitude e o acoplamento base--manipulador ao longo da
aproximação de curta distância, mas representa uma limitação do modelo
quando se deseja avaliar, em detalhe, a influência dos deslocamentos
internos de massa sobre a órbita da espaçonave.
